\chapter*{前言}
\addcontentsline{toc}{chapter}{前言}

第三册原理卷讲 AI 推理系统的机器事实，实践卷把任务拆成阶段作业。本源码卷只处理一件事：把贯穿项目 \filepath{books/cpu-volume-3/labs/linux\_cpu\_inference} 以及相关工具脚本单独成册，让读者能从构建入口一路读到测试和报告，而不是在原理正文或实践任务后面翻一段过长的源码附录。

\begin{keyidea}
源码卷不是“代码仓库的打印版”。每章先说明该组文件承担的工程责任，再给完整源码。阅读时要把函数签名、错误边界、测试断言和报告字段连起来看：一个推理引擎是否可信，不取决于某个单独函数看起来像模型，而取决于资产能校验、shape 能拒绝错误、kernel 能对照、trace 能定位、benchmark 能说明口径、CTest 能长期守住回归。
\end{keyidea}

本卷收录的主线代码是 LCQI。它从 tiny MLP 教学闭环开始，逐步覆盖 tokenizer、safetensors manifest gate、int8 kernel、packed/AVX2 路径、reference decoder、KV cache trace、GPT-2 reference path、7B compute ledger、serving SLO probe、真实 small 开源模型 smoke、自研 GPT-2 smoke 和 benchmark compare。它不是生产级推理引擎，但每个能力都要求有源码、命令、测试或报告证据。

阅读顺序建议如下。第一章先看构建入口和模型资产，明确哪些目标一定构建、哪些对标依赖外部库。第二章读公共头文件，先把数据结构和函数合同固定下来。第三章读 tiny MLP、int8 kernel 和基础 CLI，建立最小推理闭环。第四章读 reference decoder、safetensors、tokenizer 和 GPT-2 代码，理解真实模型路径为什么比模板 demo 多出许多边界检查。第五章读工具脚本、benchmark、trace、ledger、serving probe 和测试，把代码和可复查证据连起来。

源码卷的 LaTeX 文件通过 \code{lstinputlisting} 直接读取第三册工程源码。也就是说，源码不在书中复制第二份；代码更新后重新构建本卷，PDF 会读取当前仓库里的真实文件。
